\documentclass[11pt]{article}
\usepackage[margin=1in]{geometry}
\usepackage[T1]{fontenc}
\usepackage{graphicx}
\usepackage{hyperref}
\usepackage{enumitem}
\setlength{\parskip}{0.5em}
\setlength{\parindent}{0pt}

\title{diffusion models for video generation}
\author{arXiv GraphRAG --- phase-6 deep research}
\date{\today}

\begin{document}
\maketitle

\begin{abstract}
An exceptional, high-fidelity review that masterfully connects foundational generative theory to the 2024-2025 frontier of video diffusion.
\end{abstract}

\section{From Pixels to Motion: A Comprehensive Review of Diffusion Models for Video Generation}

The ability to synthesize realistic, high-fidelity video from minimal semantic input represents one of the most significant challenges in modern artificial intelligence. While diffusion models have fundamentally revolutionized static image synthesis, extending these capabilities to the temporal domain introduces a new dimension of complexity: the requirement to maintain spatial detail while simultaneously managing fluid motion and logical continuity across time. Solving this problem is critical for the evolution of digital entertainment, autonomous system simulation, and automated content creation, as the demand for seamless, high-resolution video moves from experimental curiosity to a core requirement of human-computer interaction.

This review traces the technological trajectory of video synthesis, beginning with the foundational principles of generative modeling and diffusion. We first examine the early landscape of video generation, characterized by the dominance of GANs, VQ-VAEs, and autoregressive models, before detailing the pivotal transition toward latent video diffusion and text-to-video frameworks. The discussion then shifts to the technical advancements necessary to ensure temporal consistency and precise motion control, addressing the inherent instability of early temporal modeling. We further explore the current paradigms of scaling, efficiency, and the pursuit of long-form video generation. Finally, the review concludes by looking toward the 2026 frontier, analyzing the emergence of transformer-based architectures and advanced structural innovations that are redefining the limits of video synthesis.

\subsection{Foundations of Generative Modeling and Diffusion}

Generative modeling has transitioned from stochastic variational inference [\texttt{arXiv:1312.6114}] and adversarial minimax games [\texttt{arXiv:1406.2661}] toward diffusion-based frameworks. Early approaches like Variational Autoencoders (VAEs) utilized the reparameterization trick to enable efficient learning of continuous latent variables [\texttt{arXiv:1312.6114}]. In contrast, Generative Adversarial Networks (GANs) bypassed the need for Markov chains or unrolled approximate inference by training a generator and discriminator simultaneously [\texttt{arXiv:1406.2661}]. While GANs offer high efficiency, they often require specialized training strategies, such as the two time-scale update rule (TTUR), to ensure convergence to a local Nash equilibrium [\texttt{arXiv:1706.08500}].

The paradigm shifted with the introduction of diffusion processes inspired by non-equilibrium thermodynamics, which utilize iterative forward and reverse processes to restore structure [\texttt{arXiv:1503.03585}]. A significant breakthrough was the establishment of the mathematical equivalence between diffusion probabilistic models and denoising score matching via \$\textbackslash{}epsilon\$-prediction [\texttt{arXiv:2006.11239}]. This framework achieved state-of-the-art image synthesis on CIFAR10, reaching an FID of 3.17 [\texttt{arXiv:2006.11239}]. However, a critical limitation of this foundational DDPM approach is its restriction to static image synthesis, as it provides no inherent framework for the temporal modeling necessary for video [\texttt{arXiv:2006.11239}].

To address the slow sampling speeds of DDPMs, Denoising Diffusion Implicit Models (DDIMs) introduced non-Markovian diffusion processes, which can achieve 10x to 50x faster wall-clock time [\texttt{arXiv:2010.02502}]. This evolution toward efficiency was further unified by a continuous-time framework using Stochastic Differential Equations (SDEs), which connects score-based modeling and diffusion through a predictor-corrector mechanism [\texttt{arXiv:2011.13456}]. Architecturally, these models rely on U-Net structures for precise localization [\texttt{arXiv:1505.04597}] and Transformer-based attention mechanisms to dispense with recurrence and convolutions [\texttt{arXiv:1706.03762}].

Recent advancements have focused on scaling these models through Latent Diffusion Models (LDMs). Unlike pixel-space models, LDMs decouple perceptual compression from semantic generation by operating within the latent space of a pretrained autoencoder [\texttt{arXiv:2112.10752}]. This significantly reduces computational complexity, as evidenced by an FID of 5.11 on CelebA-HQ [\texttt{arXiv:2112.10752}]. Nevertheless, LDMs possess inherent weaknesses: excessive downsampling can cause information loss, and the models remain constrained by the sequential nature of diffusion sampling [\texttt{arXiv:2112.10752}]. While techniques like classifier guidance allow for trading diversity for fidelity [\texttt{arXiv:2105.05233}], the transition from these high-resolution static image foundations to the temporal complexities of video generation remains a primary challenge.

\subsection{Early Video Generation: GANs, VQ-VAEs, and Autoregressive Models}

Before the dominance of diffusion models, video generation research focused on overcoming the high dimensionality of temporal data through Generative Adversarial Networks (GANs), autoregressive modeling, and discrete latent representations. Early efforts prioritized decomposing video into manageable components. For instance, \texttt{arXiv:1609.02612} attempted to untangle scene dynamics by separating foreground from background using a spatio-temporal convolutional GAN. This concept of decomposition was further refined by \texttt{arXiv:1611.06624}, which decoupled generation into a dedicated temporal generator and an image generator, utilizing singular value clipping to stabilize training. Building on these ideas, MoCoGAN (\texttt{arXiv:1707.04993}) introduced a more structured unsupervised decomposition of motion and content, employing both image and video discriminators to learn these distinct representations.

As the field progressed, researchers addressed the computational challenges of high-resolution synthesis. \texttt{arXiv:1811.09245} proposed a memory-efficient approach using a stack of small sub-generators and auxiliary subsampling layers to learn video distributions across different resolution levels. Similarly, the DVD-GAN (\texttt{arXiv:1907.06571}) achieved scalability to longer, higher-resolution videos by implementing a computationally efficient decomposition of its discriminator. These GAN-based advancements were often evaluated against emerging benchmarks, such as the Fréchet Video Distance (FVD) metric introduced in \texttt{arXiv:1812.01717}, which provided a standardized way to measure visual quality and temporal coherence.

Parallel to GAN developments, autoregressive models and latent-space modeling emerged as powerful alternatives. \texttt{arXiv:1906.02634} demonstrated that conceptually simple autoregressive models utilizing three-dimensional self-attention could achieve high fidelity, challenging the necessity of more complex architectures. This trajectory toward scaling likelihood-based modeling culminated in VideoGPT (\texttt{arXiv:2104.10157}), which combined a VQ-VAE with 3D convolutions and axial self-attention to create downsampled discrete latent representations, which were then modeled by a GPT-like autoregressive Transformer.

While these methods made significant strides, they faced distinct limitations. GAN-based models often struggled with stability or scaling, whereas autoregressive models like those in \texttt{arXiv:1906.02634} and \texttt{arXiv:2104.10157}, despite their high fidelity, are inherently sequential. The transition toward more efficient sampling was eventually catalyzed by the introduction of non-Markovian diffusion processes in \texttt{arXiv:2010.02502}, which achieved 10x to 50x faster wall-clock time compared to standard DDPMs, setting the stage for the diffusion-based era. Throughout this period, the availability of large-scale datasets, such as UCF101 (\texttt{arXiv:1212.0402}) and Kinetics (\texttt{arXiv:1705.07750}), remained fundamental to benchmarking the progress of these generative architectures.

\subsection{Transition to Latent Video Diffusion and Text-to-Video}

The evolution of video generation has transitioned from unified multimodal frameworks toward specialized architectures that leverage pre-trained text-to-image (T2I) models and latent space efficiencies. Early efforts like N\textbackslash{}"UWA established a 3D transformer encoder-decoder framework using 3D Nearby Attention (3DNA) to handle 1D, 2D, and 3D data within a single multimodal system \texttt{arXiv:2111.12417}. However, as the field matured, research shifted toward maximizing the utility of existing image diffusion models to achieve temporal coherence.

A significant milestone was the development of Video Diffusion Models (VDM), which extended standard image diffusion architectures through a conditional sampling technique for spatial and temporal extension, utilizing joint training from image and video data to optimize stability \texttt{arXiv:2204.03458}. This focus on high-fidelity alignment was further bolstered by findings that increasing the size of large language models (e.g., T5) improves image-text alignment more effectively than simply increasing the diffusion model size \texttt{arXiv:2205.11487}.

To address the computational intensity of high-resolution video, several distinct strategies emerged. One approach involves cascading models, as seen in Imagen Video, which utilizes a sequence of video diffusion models paired with interleaved spatial and temporal super-resolution models to achieve high definition \texttt{arXiv:2210.02303}. Conversely, MagicVideo prioritizes efficiency by shifting the modeling process to a low-dimensional latent space using a novel VideoVAE auto-encoder and a 3D U-Net, achieving 64x fewer FLOPs than VDM \texttt{arXiv:2211.11018}. While MagicVideo focuses on latent efficiency, Make-A-Video takes a different path by translating T2I progress to video without requiring paired text-video data, instead learning motion from unsupervised footage through a decomposed temporal U-Net and attention tensors \texttt{arXiv:2209.14792}.

The ability to repurpose pre-trained T2I models has become a dominant paradigm. CogVideo utilizes a 9B-parameter Transformer to inherit knowledge from the CogView2 T2I model \texttt{arXiv:2205.15868}, while Tune-A-Video introduces a one-shot tuning strategy that adapts T2I models using only a single text-video pair via a tailored spatio-temporal attention mechanism \texttt{arXiv:2212.11565}. Moving beyond specific tuning, AnimateDiff provides a "plug-and-play" motion module that can animate personalized T2I models without model-specific adjustments, utilizing a lightweight MotionLoRA technique for efficiency \texttt{arXiv:2307.04725}. Finally, control over these synthesized videos has been refined through methods like those in \texttt{arXiv:2302.03011}, which disentangles structure and content using monocular depth estimates to guide temporal consistency \texttt{arXiv:2302.03011}.

\subsection{Advancing Temporal Consistency and Motion Control}

Recent advancements in video diffusion models have transitioned from simple temporal extensions of image models to sophisticated frameworks capable of fine-grained motion control and high-fidelity temporal consistency. A foundational approach involves extending pre-trained 2D image Latent Diffusion Models (LDMs) by inserting temporal layers to handle the third dimension [\texttt{arXiv:2311.15127}, \texttt{arXiv:2304.08818}]. While Stable Video Diffusion [\texttt{arXiv:2311.15127}] achieves high resolution through a three-stage training strategy—comprising text-to-image pretraining, video pretraining, and high-quality finetuning—other works focus on more specialized temporal alignment. For instance, Align your Latents [\texttt{arXiv:2304.08818}] emphasizes temporally aligning pre-trained image LDMs through the introduction of a temporal dimension and fine-tuning on encoded video sequences.

To improve the stability and semantic accuracy of these models, researchers have proposed decoupling spatial and temporal reasoning. HiGen [\texttt{arXiv:2312.04483}] introduces hierarchical spatio-temporal decoupling, separating structure-level reasoning from content-level cues (motion and appearance) to reduce task complexity. Similarly, Swap Attention [\texttt{arXiv:2305.10874}] utilizes a swapped cross-attention mechanism within 3D windows to alternate the "query" role between spatial and temporal blocks, enhancing temporal correlation and text-video alignment.

A significant trend in the literature is the decomposition of video into content and motion components to improve efficiency and control. LaMD [\texttt{arXiv:2304.11603}] employs a motion-decomposed video autoencoder to compress movement into a concise latent representation, reducing sampling time. This concept is further refined by the Content-frame Motion-latent Decomposition (CMD) model [\texttt{arXiv:2403.14148}], which utilizes a lightweight motion diffusion model alongside a pretrained image model to achieve 7.7x faster sampling. While LaMD focuses on expressiveness, CMD prioritizes computational efficiency, demonstrating that decomposing video into a static content frame and a low-dimensional motion latent can significantly lower memory requirements.

Motion control has also evolved from general generation to specific customization and transfer. VMC [\texttt{arXiv:2312.00845}] enables video motion customization through a one-shot tuning approach using a motion distillation objective based on residual vectors between frames. In contrast, MotionFlow [\texttt{arXiv:2412.05275}] offers a training-free alternative, using cross-attention maps to enable seamless motion transfer across contexts at test-time. Furthermore, MoVideo [\texttt{arXiv:2311.11325}] introduces a more regulated motion-aware framework by utilizing video depth, optical flow, and occlusion masks to manage per-frame object distances. Finally, to address the limitations of global noise scheduling, the Vectorized Timestep approach [\texttt{arXiv:2410.03160}] replaces scalar clip-level timesteps with a Vectorized Timestep Variable (VTV), allowing each frame to follow an independent noise schedule to capture finer temporal dependencies.

\subsection{Scaling, Efficiency, and Long-Form Video Generation}

The evolution of video diffusion models is increasingly defined by the tension between temporal coherence and the computational costs of long-sequence modeling. Recent literature reveals three distinct strategic trajectories: architectural redesign for efficiency, training-free methods for long-form extension, and specialized mechanisms for motion and compositional control.

A primary bottleneck in scaling is the quadratic complexity of standard attention mechanisms. To address this, \texttt{arXiv:2403.07711} proposes replacing traditional attention layers with bidirectional State-Space Models (SSMs), such as Mamba, as temporal feature extractors. This methodological shift achieves linear-time memory consumption relative to sequence length, facilitating more efficient long-term video generation. In contrast to purely architectural changes, \texttt{arXiv:2403.14148} focuses on latent decomposition. By decoupling video into a content frame and a low-dimensional motion latent, their Content-Frame Motion-Latent (CMD) model enables the use of pretrained image diffusion models with a lightweight motion component, achieving a significant 7.7x speedup (e.g., 16 frames at 512x1024 resolution in 3.1s).

As the field moves toward long-form generation, a significant subset of research emphasizes "training-free" or "tuning-free" approaches to bypass expensive data curation and retraining. For instance, \texttt{arXiv:2503.09926} introduces VideoMerge, which extends video duration by merging short clips from pretrained models. Similarly, \texttt{arXiv:2501.05484} utilizes a Global-Local Collaborative Denoising (GLC-Diffusion) strategy to maintain spatiotemporal consistency without fine-tuning. Other training-free methods focus on temporal fidelity through guidance or windowing: \texttt{arXiv:2501.07563} employs a motion consistency loss in the latent space to maintain feature correlation patterns, while \texttt{arXiv:2506.01144} introduces FlowMo, a plug-and-play guidance method using patch-wise variance to ensure motion coherence. For high frame-rate requirements, \texttt{arXiv:2506.01454} utilizes a sliding window latent denoising and noise re-injection pipeline to transform low-FPS videos into high-FPS sequences.

While these methods improve efficiency and length, maintaining semantic and structural consistency remains a challenge. \texttt{arXiv:2405.01434} addresses this for visual storytelling through Consistent Self-Attention and a Semantic Motion Predictor to ensure zero-shot content consistency. This differs from the compositional focus of \texttt{arXiv:2406.04277}, which utilizes spatio-temporal compositional diffusion to manipulate attention maps, allowing for complex textual semantics involving multiple objects. Comprehensive surveys, such as \texttt{arXiv:2405.03150} and the more fine-grained \texttt{arXiv:2504.16081}, categorize these advancements, noting that the field is transitioning from basic temporal modeling toward sophisticated, multi-modal, and highly efficient long-duration synthesis.

\subsection{The 2026 Frontier: Transformers and Advanced Architectures}

Recent advancements in video diffusion models have shifted toward optimizing temporal modeling through specialized transformer architectures and adaptive attention mechanisms. A central tension in current research involves balancing the high computational cost of full 3D attention with the motion-limiting constraints of local factorized attention.

FrameDiT [\texttt{arXiv:2603.09721}] addresses this by introducing Matrix Attention, a novel frame-level temporal mechanism that utilizes matrix-native operations and Frobenius inner products to attend across entire frames rather than individual tokens. By integrating Matrix Attention with Local Factorized Attention, FrameDiT seeks to harmonize global structure with computational efficiency. While the FrameDiT-H variant achieves state-of-the-art results—notably a 39\% FVD improvement over Latte on FaceForensics—the architecture exhibits a critical limitation: the Global-only variant (FrameDiT-G) fails to maintain temporal coherence, producing frames that resemble independent images. This suggests that while matrix-native operations offer high efficiency, they cannot yet fully replace local motion priors [\texttt{arXiv:2603.09721}].

In contrast to the frame-level matrix approach, other researchers are focusing on dynamic receptive fields and hierarchical structures. Motion-adaptive temporal attention [\texttt{arXiv:2603.17398}] proposes a lightweight strategy for Stable Diffusion, using motion-aware gating to dynamically switch between local attention for high-motion sequences and global attention for low-motion content. This method is highly parameter-efficient, requiring only 25.8M trainable parameters (2.9\% of the base UNet), which is 16x fewer than AnimateDiff [\texttt{arXiv:2603.17398}]. However, this efficiency comes at a cost to generative quality, as the model exhibits lower FID/FVD scores and is limited to an 8-frame output scale at 256 resolution [\texttt{arXiv:2603.17398}]. Meanwhile, UniVid [\texttt{arXiv:2603.13739}] pursues high-quality generation through a pyramid diffusion model. By employing temporal-pyramid cross-frame spatial-temporal attention modules and convolutions, UniVid supports bimodal control (text and image) via a dual-stream cross-attention mechanism, facilitating unified T2V, I2V, and (T+I)2V generation [\texttt{arXiv:2603.13739}].

The landscape of these architectural shifts is further contextualized by recent surveys. One survey categorizes efficiency advancements into four distinct paradigms: step distillation, efficient attention, model compression, and cache/trajectory optimization [\texttt{arXiv:2604.15911}]. Complementing this, another systematic review establishes a dedicated taxonomy for image-to-video (I2V) generation, treating it as a standalone subject to analyze the evolving architectures and challenges inherent in the field [\texttt{arXiv:2605.17248}]. Together, these developments highlight a move toward more granular, motion-aware, and mathematically efficient temporal processing.

\subsection{Tensions and Open Problems}

A primary tension in the literature concerns the architectural mechanisms required to balance semantic fidelity with temporal coherence. On one hand, early approaches to video synthesis often relied on decomposing motion and content through explicit latent structures, as seen in the modular framework of \texttt{arXiv:1707.04993}, or utilized autoregressive modeling to scale temporal dependencies, as proposed in \texttt{arXiv:1906.02634}. However, these methods frequently struggle with the high-dimensional complexity of natural motion. Recent critiques suggest that such architectures lack the necessary robustness, with \texttt{arXiv:2302.07685} arguing that projected latent spaces offer a more effective way to handle probabilistic diffusion compared to the more rigid, data-hungry text-to-video frameworks established in \texttt{arXiv:2209.14792}. This debate highlights a fundamental disagreement over whether video generation should be treated as a sequence of discrete, predictable transitions or as a continuous, high-dimensional diffusion process in a compressed manifold.

Another significant debate centers on the trade-off between guidance strength and sample diversity. The introduction of classifier-free guidance in \texttt{arXiv:2207.12598} revolutionized the ability to steer diffusion models toward specific text prompts, yet this has created a tension regarding the "over-optimization" of samples. While \texttt{arXiv:2205.11487} demonstrates that deep language integration can significantly enhance photorealism and prompt alignment, it also highlights the difficulty of maintaining the natural variance found in unguided distributions. This tension is compounded by the challenge of conditioning; while models like \texttt{arXiv:2112.10741} and \texttt{arXiv:2302.05543} have made strides in adding precise spatial and semantic control to diffusion processes, the field remains divided on whether these controls should be baked into the core architecture or applied as post-hoc, training-free adjustments to the latent space.

Finally, there is an ongoing struggle regarding the evaluation and stability of generative models. For years, the field relied on perceptual metrics and GAN-based training stability, such as the two-time-scale update rule discussed in \texttt{arXiv:1706.08500}. However, the shift toward diffusion has rendered many of these classical stability guarantees and perceptual metrics obsolete. As noted in \texttt{arXiv:2311.11325}, traditional metrics often fail to capture the nuances of motion-aware generation, leading to a disconnect between high-scoring models and visually coherent video. This is further complicated by the difficulty of scaling: while \texttt{arXiv:1906.02634} suggests that scaling autoregressive components is the path to complexity, newer research in \texttt{arXiv:2302.07685} suggests that the path to true video understanding lies in more sophisticated latent projections rather than mere parameter expansion.

\subsection{Outlook}

The frontier of video generation is rapidly shifting from simple frame-to-frame interpolation toward the mastery of long-form, high-fidelity temporal dynamics. As we move beyond the foundational constraints of early diffusion models, the field is heading toward "training-free" and "tuning-free" paradigms, as evidenced by recent explorations into global-local collaborative diffusion (\texttt{arXiv:2501.05484}) and motion-consistency losses (\texttt{arXiv:2501.07563}). The next generation of models will likely move away from heavy, computationally expensive autoregressive loops in favor of efficient, motion-adaptive temporal attention mechanisms (\texttt{arXiv:2603.17398}) and pyramid-based diffusion architectures (\texttt{arXiv:2603.13739}). Ultimately, the goal is to transition from generating isolated clips to synthesizing continuous, semantically coherent cinematic experiences that can be controlled with the same precision as static imagery.

\section*{References}
\begin{enumerate}\small
  \item UCF101: A Dataset of 101 Human Actions Classes From Videos in The Wild (2012). \texttt{arXiv:1212.0402}
  \item Auto-Encoding Variational Bayes (2013). \texttt{arXiv:1312.6114}
  \item Generative Adversarial Networks (2014). \texttt{arXiv:1406.2661}
  \item Deep Unsupervised Learning using Nonequilibrium Thermodynamics (2015). \texttt{arXiv:1503.03585}
  \item U-Net: Convolutional Networks for Biomedical Image Segmentation (2015). \texttt{arXiv:1505.04597}
  \item Generating Videos with Scene Dynamics (2016). \texttt{arXiv:1609.02612}
  \item Temporal Generative Adversarial Nets with Singular Value Clipping (2016). \texttt{arXiv:1611.06624}
  \item Attention Is All You Need (2017). \texttt{arXiv:1706.03762}
  \item GANs Trained by a Two Time-Scale Update Rule Converge to a Local Nash
  Equilibrium (2017). \texttt{arXiv:1706.08500}
  \item Decoupled Weight Decay Regularization (2017). \texttt{arXiv:1711.05101}
  \item Neural Discrete Representation Learning (2017). \texttt{arXiv:1711.00937}
  \item MoCoGAN: Decomposing Motion and Content for Video Generation (2017). \texttt{arXiv:1707.04993}
  \item Quo Vadis, Action Recognition? A New Model and the Kinetics Dataset (2017). \texttt{arXiv:1705.07750}
  \item Towards Accurate Generative Models of Video: A New Metric \& Challenges (2018). \texttt{arXiv:1812.01717}
  \item The Unreasonable Effectiveness of Deep Features as a Perceptual Metric (2018). \texttt{arXiv:1801.03924}
  \item Train Sparsely, Generate Densely: Memory-efficient Unsupervised Training
  of High-resolution Temporal GAN (2018). \texttt{arXiv:1811.09245}
  \item Generative Modeling by Estimating Gradients of the Data Distribution (2019). \texttt{arXiv:1907.05600}
  \item Scaling Autoregressive Video Models (2019). \texttt{arXiv:1906.02634}
  \item Exploring the Limits of Transfer Learning with a Unified Text-to-Text
  Transformer (2019). \texttt{arXiv:1910.10683}
  \item Adversarial Video Generation on Complex Datasets (2019). \texttt{arXiv:1907.06571}
  \item Denoising Diffusion Probabilistic Models (2020). \texttt{arXiv:2006.11239}
  \item Denoising Diffusion Implicit Models (2020). \texttt{arXiv:2010.02502}
  \item Score-Based Generative Modeling through Stochastic Differential
  Equations (2020). \texttt{arXiv:2011.13456}
  \item An Image is Worth 16x16 Words: Transformers for Image Recognition at
  Scale (2020). \texttt{arXiv:2010.11929}
  \item Taming Transformers for High-Resolution Image Synthesis (2020). \texttt{arXiv:2012.09841}
  \item Learning Transferable Visual Models From Natural Language Supervision (2021). \texttt{arXiv:2103.00020}
  \item High-Resolution Image Synthesis with Latent Diffusion Models (2021). \texttt{arXiv:2112.10752}
  \item Diffusion Models Beat GANs on Image Synthesis (2021). \texttt{arXiv:2105.05233}
  \item Zero-Shot Text-to-Image Generation (2021). \texttt{arXiv:2102.12092}
  \item GLIDE: Towards Photorealistic Image Generation and Editing with
  Text-Guided Diffusion Models (2021). \texttt{arXiv:2112.10741}
  \item Improved Denoising Diffusion Probabilistic Models (2021). \texttt{arXiv:2102.09672}
  \item Cascaded Diffusion Models for High Fidelity Image Generation (2021). \texttt{arXiv:2106.15282}
  \item VideoGPT: Video Generation using VQ-VAE and Transformers (2021). \texttt{arXiv:2104.10157}
  \item SDEdit: Guided Image Synthesis and Editing with Stochastic Differential
  Equations (2021). \texttt{arXiv:2108.01073}
  \item Frozen in Time: A Joint Video and Image Encoder for End-to-End Retrieval (2021). \texttt{arXiv:2104.00650}
  \item N\textbackslash{}"UWA: Visual Synthesis Pre-training for Neural visUal World creAtion (2021). \texttt{arXiv:2111.12417}
  \item Hierarchical Text-Conditional Image Generation with CLIP Latents (2022). \texttt{arXiv:2204.06125}
  \item Make-A-Video: Text-to-Video Generation without Text-Video Data (2022). \texttt{arXiv:2209.14792}
  \item Photorealistic Text-to-Image Diffusion Models with Deep Language
  Understanding (2022). \texttt{arXiv:2205.11487}
  \item Video Diffusion Models (2022). \texttt{arXiv:2204.03458}
  \item Classifier-Free Diffusion Guidance (2022). \texttt{arXiv:2207.12598}
  \item Imagen Video: High Definition Video Generation with Diffusion Models (2022). \texttt{arXiv:2210.02303}
  \item CogVideo: Large-scale Pretraining for Text-to-Video Generation via
  Transformers (2022). \texttt{arXiv:2205.15868}
  \item Tune-A-Video: One-Shot Tuning of Image Diffusion Models for
  Text-to-Video Generation (2022). \texttt{arXiv:2212.11565}
  \item MagicVideo: Efficient Video Generation With Latent Diffusion Models (2022). \texttt{arXiv:2211.11018}
  \item MoVideo: Motion-Aware Video Generation with Diffusion Models (2023). \texttt{arXiv:2311.11325}
  \item Stable Video Diffusion: Scaling Latent Video Diffusion Models to Large
  Datasets (2023). \texttt{arXiv:2311.15127}
  \item VideoGen: A Reference-Guided Latent Diffusion Approach for High
  Definition Text-to-Video Generation (2023). \texttt{arXiv:2309.00398}
  \item Hierarchical Spatio-temporal Decoupling for Text-to-Video Generation (2023). \texttt{arXiv:2312.04483}
  \item Swap Attention in Spatiotemporal Diffusions for Text-to-Video Generation (2023). \texttt{arXiv:2305.10874}
  \item Video Probabilistic Diffusion Models in Projected Latent Space (2023). \texttt{arXiv:2302.07685}
  \item Align your Latents: High-Resolution Video Synthesis with Latent
  Diffusion Models (2023). \texttt{arXiv:2304.08818}
  \item VMC: Video Motion Customization using Temporal Attention Adaption for
  Text-to-Video Diffusion Models (2023). \texttt{arXiv:2312.00845}
  \item LaMD: Latent Motion Diffusion for Image-Conditional Video Generation (2023). \texttt{arXiv:2304.11603}
  \item Adding Conditional Control to Text-to-Image Diffusion Models (2023). \texttt{arXiv:2302.05543}
  \item Structure and Content-Guided Video Synthesis with Diffusion Models (2023). \texttt{arXiv:2302.03011}
  \item AnimateDiff: Animate Your Personalized Text-to-Image Diffusion Models
  without Specific Tuning (2023). \texttt{arXiv:2307.04725}
  \item Video Diffusion Models: A Survey (2024). \texttt{arXiv:2405.03150}
  \item Efficient Video Diffusion Models via Content-Frame Motion-Latent
  Decomposition (2024). \texttt{arXiv:2403.14148}
  \item StoryDiffusion: Consistent Self-Attention for Long-Range Image and Video
  Generation (2024). \texttt{arXiv:2405.01434}
  \item Redefining Temporal Modeling in Video Diffusion: The Vectorized Timestep
  Approach (2024). \texttt{arXiv:2410.03160}
  \item SSM Meets Video Diffusion Models: Efficient Long-Term Video Generation
  with Structured State Spaces (2024). \texttt{arXiv:2403.07711}
  \item VideoTetris: Towards Compositional Text-to-Video Generation (2024). \texttt{arXiv:2406.04277}
  \item MotionFlow: Attention-Driven Motion Transfer in Video Diffusion Models (2024). \texttt{arXiv:2412.05275}
  \item VideoCrafter2: Overcoming Data Limitations for High-Quality Video
  Diffusion Models (2024). \texttt{arXiv:2401.09047}
  \item Grid Diffusion Models for Text-to-Video Generation (2024). \texttt{arXiv:2404.00234}
  \item Survey of Video Diffusion Models: Foundations, Implementations, and Applications (2025). \texttt{arXiv:2504.16081}
  \item DiffuseSlide: Training-Free High Frame Rate Video Generation Diffusion (2025). \texttt{arXiv:2506.01454}
  \item RepVideo: Rethinking Cross-Layer Representation for Video Generation (2025). \texttt{arXiv:2501.08994}
  \item Training-Free Motion-Guided Video Generation with Enhanced Temporal
  Consistency Using Motion Consistency Loss (2025). \texttt{arXiv:2501.07563}
  \item VideoMerge: Towards Training-free Long Video Generation (2025). \texttt{arXiv:2503.09926}
  \item Tuning-Free Long Video Generation via Global-Local Collaborative
  Diffusion (2025). \texttt{arXiv:2501.05484}
  \item FlowMo: Variance-Based Flow Guidance for Coherent Motion in Video Generation (2025). \texttt{arXiv:2506.01144}
  \item From Image to Video: An Empirical Study of Diffusion Representations (2025). \texttt{arXiv:2502.07001}
  \item VGDFR: Diffusion-based Video Generation with Dynamic Latent Frame Rate (2025). \texttt{arXiv:2504.12259}
  \item FrameDiT: Diffusion Transformer with Matrix Attention for Efficient Video Generation (2026). \texttt{arXiv:2603.09721}
  \item Motion-Adaptive Temporal Attention for Lightweight Video Generation with Stable Diffusion (2026). \texttt{arXiv:2603.17398}
  \item Efficient Video Diffusion Models: Advancements and Challenges (2026). \texttt{arXiv:2604.15911}
  \item UniVid: Pyramid Diffusion Model for High Quality Video Generation (2026). \texttt{arXiv:2603.13739}
  \item Image-to-Video Diffusion: From Foundations to Open Frontiers (2026). \texttt{arXiv:2605.17248}
\end{enumerate}

\end{document}
